\section{Limitations}\label{sec:limitations}

\begin{enumerate}
\item \emph{Scope.} One task family, small recurrent networks (about 40,000
  parameters), three seeds per condition and fixed budgets: we report
  observed counts, not population frequencies, and cannot claim that the
  separations hold for other tasks, architectures or longer training.
\item \emph{Constructed tests.} The law panels and the census locate failures
  but do not estimate how often they occur under the natural distribution.
\item \emph{Confounded comparisons.} Comparisons across studies (r11 versus
  r12; anything involving r10) change more than one factor and are not
  attributed to a single cause; within r12, the version comparison changes
  supplied sums, sums accuracy and sums supervision together and does not
  isolate access.
\item \emph{Internal structure.} Probe declines show reduced measured
  readability, not removal of information, for the tested probe families and
  support; probe boards are disjoint between fitting and evaluation but not
  necessarily unseen during network training, and rare r10 sums (28--36) lie
  outside the registered scope. The connection experiment changes the route,
  not the internal state, and the failed swap selectivity excludes no other
  representation. None of these establishes how the free networks represent
  or use the sums internally.
\item \emph{Encoding.} The encoding experiment compares two encodings and
  does not establish a precision ceiling; the supplementary cutoff
  discrimination is established on that support, not across all cutoff
  boards, and does not diagnose the numerical-encoding failures.
\item \emph{Mechanisms.} Erosion depends on the prediction gradient, but the
  runs do not separate gradient direction, displacement and answer margins,
  so interference is not established; the residual failures of networks
  given exact sums are located but not explained, and a pass is not
  exhaustive correctness.
\end{enumerate}
